\Needspace{8\baselineskip}\originalchapter{冷等离子体的电磁波}\label{ch:cold}
\originalsection{介电张量的来源和约定}
取实场为复振幅实部, 时间因子 $\e^{-\ii\omega t}$, $\mathbf B_0=B_0\mathbf e_z$. 定义带符号 $\Omega_s=q_sB_0/m_s$, 正数 $\omega_{ps}^2=n_sq_s^2/(\epsilon_0m_s)$. 冷流体扰动满足 $-\ii\omega m_s\mathbf v_s=q_s(\mathbf E+\mathbf v_s\times\mathbf B_0)$. 解二维线性方程并求 $\mathbf J=\sum n_sq_s\mathbf v_s=-\ii\omega\epsilon_0(\boldsymbol\epsilon-I)\mathbf E$, 得
\begin{equation}
 \boldsymbol\epsilon=\begin{pmatrix}S&-\ii D&0\\\ii D&S&0\\0&0&P\end{pmatrix},\quad
 S=1-\sum_s\frac{\omega_{ps}^2}{\omega^2-\Omega_s^2},\quad
 D=\sum_s\frac{\Omega_s\omega_{ps}^2}{\omega(\omega^2-\Omega_s^2)},\quad
 P=1-\sum_s\frac{\omega_{ps}^2}{\omega^2}.
\end{equation}
此处电子 $\Omega_e<0$. $R=S+D=1-\sum\omega_{ps}^2/[\omega(\omega+\Omega_s)]$, $L=S-D$; 电子 R 支在正 $\omega=|\Omega_e|$ 有回旋共振. 若把 $\Omega_e$ 改定义为正大小却照抄 D 符号, 极化会反向.

\originalsection{传播方向决定分支}
Maxwell 方程消去 B 得 $[\boldsymbol\epsilon-n_r^2I+\mathbf n_r\mathbf n_r]\mathbf E=0$, $\mathbf n_r=c\mathbf k/\omega$. 平行传播有 $n_r^2=R,L$, 横场圆偏振, 另有 P=0 的纵向电子波. 垂直传播普通支 $n_r^2=P$、E沿B; 非常支 $n_r^2=RL/S$, E在垂直B平面且一般含纵向分量. 任意夹角不能用这四个端点分支替代, 应解完整行列式.

\originalsection{任意夹角的两个电磁根}
为落实22.611J 2003讲义第5章的任意角计算，令$\mathbf k=k(\sin\theta,0,\cos\theta)$，$N^2=c^2k^2/\omega^2$。把上一节张量代入Maxwell方程，三个电场分量满足
\begin{equation}
\begin{pmatrix}
S-N^2\cos^2\theta&-\ii D&N^2\sin\theta\cos\theta\\
\ii D&S-N^2&0\\
N^2\sin\theta\cos\theta&0&P-N^2\sin^2\theta
\end{pmatrix}
\begin{pmatrix}E_x\\E_y\\E_z\end{pmatrix}=0 .
\end{equation}
展开行列式，纯介质项$C=PRL$，$N^2$与$N^4$项分别给
\begin{equation}
A N^4-B N^2+C=0,\qquad
A=S\sin^2\theta+P\cos^2\theta,\quad
B=RL\sin^2\theta+PS(1+\cos^2\theta).
\end{equation}
因$R=S+D,L=S-D$，判别式可化为
\begin{equation}
B^2-4AC=(RL-PS)^2\sin^4\theta+4P^2D^2\cos^2\theta\equiv F^2,
\qquad N_\pm^2=\frac{B\pm F}{2A}.
\end{equation}
$A=0$时要回到原多项式判断有限根及趋于无穷的共振支，不能直接除零。给定有限根，第二、三行还决定偏振：
$E_y=-\ii D E_x/(S-N^2)$、
$E_z=-N^2\sin\theta\cos\theta\,E_x/(P-N^2\sin^2\theta)$，分母为零的特殊点改用原矩阵零空间，不机械套比值。

端点检验是推导的一部分。$\theta=0$时四次式分解为$P(N^2-R)(N^2-L)=0$，两横向圆偏振支是$R,L$；$\theta=\pi/2$时为$S(N^2-P)(N^2-RL/S)=0$，回到O/X两支。举一个明确的无碰撞电子响应代数例子：取$|\Omega_e|=1,\omega=2,\omega_{pe}^2=1$，并忽略离子高频响应，得到$S=2/3,D=-1/6,P=3/4,R=1/2,L=5/6$。在$\theta=\pi/4$，$A=17/24,B=23/24,C=5/16$，故$N_\pm^2=(23\pm\sqrt{19})/34\simeq0.80467,0.54827$；它们位于两组端点之间，且对应不同的椭圆偏振。角度变化时应连续跟踪根与偏振，不能只按“加号就是某端点名称”作全区命名。

再问何时出现静电主导的高$N$支。对矩阵沿$\hat{\mathbf k}$投影，最高阶电磁项抵消，必要的纵向介电条件为$\hat{\mathbf k}\cdot\boldsymbol\epsilon\cdot\hat{\mathbf k}=A=0$。它是$N^2\to\infty$的极限条件，并不意味着某个有限$N$下$E_\perp$严格为零：有限波长的剩余横向方程通常要求$E_\perp/E_\parallel=O(N^{-2})$。直接把$E_\perp=0$代回三个分量式可能留下非零横向介质耦合。这个区别决定混合模在共振附近是否仍可耦合外来电磁波。

在无碰撞模型中, $n_r^2=0$ 为截止, $n_r^2\to\infty$ 为共振. 负 $n_r^2$ 是倏逝, 不应取其对数后画成实传播曲线. 因为群速度和能流需包含色散介质储能, 折射率很大并不等于信号速度无穷大. 在理想共振附近冷模型常需温度、碰撞或非局域修正.

\originalsection{截止层的局部结构}
未磁化分层等离子体中, 横波有 $E''+k_0^2[1-\omega_p^2(z)/\omega^2]E=0$. 近截止 $n(z)=n_c(1+z/a)$, 化为 $E''-k_0^2zE/a=0$. 令 $\zeta=(k_0^2/a)^{1/3}z$, 解为 Ai、Bi; 高密度一侧物理解为衰减 Ai. 低密度侧 Ai 的正弦渐近式包含幅度相同的入射与反射波, 因无耗散反射率为1. WKB 在截止处失效, Airy 层补上相位连接. 这比把 k=0 直接代进均匀平面波更能解释反射.

\originalsection{密度坡中的可达性}
若电子响应且 $u=x/\ell$, $w=\omega/|\Omega_e|$, 密度标度令 $\omega_p^2(n_0)=\Omega_e^2$, 则 O 截止 u=w\ensuremath{{}^{2}}. X 支
\begin{equation}
 n_X^2=\frac{(w^2-u)^2-w^2}{w^2(w^2-u-1)}.
\end{equation}
$w>1$ 时依次遇 $u_R=w(w-1)$ 截止、$u_H=w^2-1$ 共振、$u_L=w(w+1)$ 截止. 共振后的传播带虽然存在, 从真空来的波仍需穿过前面的倏逝层. 波加热设计必须同时检查极化、可达性和吸收, 仅画共振位置不足以证明能把能量送到那里.
